\chapter{Introduction}\label{chp:intro}

% starting arabic page numbering
\pagenumbering{arabic}

% prevent orphan/widow lines (1-2 lines alone at top/bottom of page)
\widowpenalty=10000
\clubpenalty=10000

Cloud computing has dramatically changed how organizations utilize their computer processing capabilities. Today's modern data centers are using clusters made up of thousands of computers to provide both interactive web-based applications and large-scale analytical applications. Machine learning (ML) applications are among the most computationally intensive workloads processed in today's data centers and often require graphics processing units (GPUs) to train models and perform inference operations. To keep GPU clusters running efficiently, they need high throughput, high hardware utilization, and low waiting time for jobs simultaneously. However, achieving all three of these goals at once is far from straightforward.

This difficulty is exemplified in the large-scale production environments of today. Organizations like Alibaba have created large, dedicated GPU clusters with thousands of accelerators that run continuously with a variety of different workload patterns, ranging from short pre-processing scripts to multi-day model training runs~\cite{weng2022mlaas}. Moreover, the resources utilized by these workloads vary greatly. For example, some workloads may fully saturate a GPU's memory and almost never touch the central processing unit (CPU), while other workloads may do the opposite. Therefore, if a scheduler were to treat all workloads identically, there would be a significant likelihood that some jobs would be starved for available resources and the capacity provided to others could go unused.

Measurements taken in the field show that this is an actual problem in current practice, rather than simply a hypothetical one. Researchers have studied the performance of large-scale shared GPU infrastructures and found that average hardware utilization levels fall significantly short of what the installed capacity would suggest. The researchers indicated that the median GPU usage was about 52\% in all production cluster environments, while even lower, about 28.26\%, for larger distributed jobs. They also stated that the most common causes of these low usage rates were strict scheduling practices, such as gang scheduling, and localities that restricted resource access~\cite{jeon2019analysis}. This is an important finding: the limiting factor is not the availability of sufficient hardware; it is the policies governing the usage of that hardware.

\clearpage
A key mechanism for making those trade-off decisions is the scheduler, which is responsible for determining when and where each submitted job should run. If a scheduler does not know ahead of time how much compute time will be required by each job, then it will make decisions based on simplistic ordering rules that either result in wasted resources or result in many short-duration jobs being blocked by longer-duration jobs. Providing the scheduler with accurate estimates of the runtime needed to complete each job via the application of ML is a potential area for improvement in cluster efficiency and is the central theme of this thesis.

\needspace{5\baselineskip}
\section{Problem Definition and Context}

Over the last ten years, the need for computing resources to support both ML and artificial intelligence (AI) has grown exponentially. Thousands of GPU hours are required to train one of today's state-of-the-art deep learning (DL) models. Additionally, when it comes to high-volume, continuous, low-latency acceleration for inference services (IFS), access to these accelerators is critical. Because of this rapid increase in demand, GPU clusters now play a primary role in today's computing environments.

Cloud computing has been able to meet increasing demand for large amounts of elastic and heterogeneous hardware in order to be used by multiple individuals and applications at one time. At production scale, there are thousands of CPUs and GPUs being combined into large cluster systems to support a variety of applications including short-term jobs running on small pieces of data as well as long-term applications such as multi-day model training. As cloud computing continues to scale, it is becoming much harder to sustainably maintain high levels of utilization from all of the available hardware while supporting an ever-increasing number of concurrent users and application types.

The high amount of operational complexity associated with running large-scale clusters has been widely described. For example, the Borg system developed by Google manages a multitude of jobs on many machines at one time. It uses techniques like task-packing, resource over-allocation, and preemptive prioritization to allow for acceptable utilization levels~\cite{verma2015large}.

\clearpage
Using data from previous workloads as input for ML models, Resource Central was shown to be able to accurately model how virtual machines (VMs) would act given those inputs, and thus use those models to help proactively place VMs in order to minimize waste~\cite{cortez2017resource}. In a study of Microsoft's multi-tenant GPU cluster, it was found that there were very low average levels of utilization of GPUs (only about 52\%) and it was determined that the poor distribution of jobs and inter-job interference were both major contributors to the low usage~\cite{jeon2019analysis}. Collectively, these studies demonstrate that static and rule-based methods of scheduling are insufficient in modern production environments.

Alibaba's GPU cluster has a similar story. Alibaba's Platform for Artificial Intelligence (PAI) runs on roughly 1,800 machines using about 6,500 GPUs, supporting various combinations of training and inference workloads~\cite{weng2022mlaas}. Alibaba found that memory contention and co-located workload unpredictability were the two biggest obstacles to higher efficiency~\cite{guo2019limits}.

High utilization rates of resources are fundamentally challenging since cloud-based workloads are inherently chaotic. For example, the Alibaba PAI trace indicates that job submissions exhibit patterns that repeat on a daily basis; however, they are also punctuated with spikes that can rapidly exceed available capacity~\cite{weng2022mlaas}. Alongside this temporal variability, job execution times within GPU clusters also exhibit extremely skewed distributions. In systems literature, this is generally referred to as the ``mice and elephants'' problem. A large number of jobs are short-lived, interactive ``mice'' that complete in minutes. Conversely, a very small percentage of jobs involve large-scale training; these are ``elephant-sized'' runs that last for many days, consuming a disproportionate amount of cluster time. As a result, this variance creates resource fragmentation problems. Resource fragmentation occurs when there is available CPU capacity at a particular node but insufficient GPU memory to allow the use of those available resources~\cite{jeon2019analysis}. The ability of a scheduler to schedule jobs appropriately and anticipate potential increases in demand is severely impaired due to the lack of knowledge about the duration of each job. Therefore, a scheduling environment is created which results in consistently unused resources and increasing queue lengths.

\clearpage
The most basic scheduling algorithm, First-In-First-Out (FIFO), schedules jobs based on when they arrived. Although equitable, a single long-running task could block many short tasks. Shortest Job First (SJF) solves this by processing the shortest jobs first, and is proven to be the best schedule for minimizing average wait times~\cite{kleinrock1975queueing}. The problem with SJF is that an estimate of how long a job will run needs to be made prior to its execution, and such estimates are almost never available in practice.

A very popular approach in high-performance computing (HPC) has been to have users provide an estimated upper limit on how long they think a job will run. These estimates are then used by the scheduler as a proxy for actual running time. Recent studies of large-scale HPC batch workloads have shown that users routinely overestimate these limits by factors ranging from two to 10, or even higher~\cite{tsafrir2007backfill}. As such, the scheduler makes scheduling decisions based upon inaccurate, noisy, and biased information, commonly backfilling the queue with likely non-conflicting jobs while simultaneously blocking numerous smaller jobs that are waiting for resources. This leads us to replace user-provided estimates with ML-based predictions using workload data.

Tetris, developed by Grandl et al., utilized multi-resource packing to minimize fragmentation~\cite{grandl2014multi}, and Quincy, developed by Isard et al., implemented network-flow optimization to find an optimal solution between fairness and locality of data~\cite{isard2009quincy}. TetriSched expanded on this concept through adaptive re-planning for heterogeneous hardware~\cite{tumanov2016tetrisched}, and Decima demonstrated the application of deep reinforcement learning (DRL) to develop scheduling policies based on real-world workload data~\cite{mao2019learning}. All of these approaches shared one common understanding in that, if a scheduler could make accurate predictions about runtime, it would be possible to simulate SJF without the need for having ``oracle'' knowledge. Gradient boosting models, such as XGBoost (XGB)~\cite{chen2016xgboost} and LightGBM (LGBM)~\cite{ke2017lightgbm}, have been shown to be successful for predicting outcomes on tabular data. Recurrent architectures, specifically Long Short-Term Memory (LSTM) networks~\cite{hochreiter1997long}, are particularly adept at identifying sequences within time.

\clearpage
In contrast to unstructured data like images or text, on which DL performs well due to structured pattern exploitation, cluster traces are tabular, where each job is represented by a fixed set of features such as GPU count, CPU request, and submission time. Gradient boosting models have recently been demonstrated to be competitive with, and often perform better than, DL networks on tabular data~\cite{shwartzziv2022tabular}. Therefore, one of the key questions addressed in this thesis is whether this advantage also applies to GPU cluster traces.

\needspace{5\baselineskip}
\section{Motivation, Definition, and Scope}

The underlying reason for undertaking this study is the fact that historically, both runtime predictions and job scheduling have been addressed as independent issues within an overarching ``Predict-then-Optimize'' (PtO) approach. Under PtO, predictive models are developed primarily to minimize error at a specific point in time, largely ignoring downstream scheduling objectives. However, predicting the duration of jobs is more than just a statistical task; these insights serve as a direct tool to drive smarter scheduling choices. Thus, this thesis supports a transition away from the conventional PtO paradigm toward a Decision-Focused Learning (DFL) framework~\cite{mandi2023decision}. Within DFL, the overall success of a predictor is judged by its capacity to optimize system-level performance measures, such as job completion times (JCTs), rather than solely focusing on raw accuracy. Isolated reports on prediction accuracy fail to show whether or how these estimates actually contribute to superior scheduling outcomes, especially considering how extreme outlier values can heavily distort a model's loss function.

By providing a full flow from raw data to the evaluation of schedules in this thesis, it closes the gap that exists for traces such as those used in this thesis with the Alibaba PAI GPU cluster trace~\cite{weng2022mlaas}. It includes all aspects of data cleaning and feature engineering, including a sweep-line-based computing method for cluster utilization. Furthermore, it includes the design and comparison of both tree-based models (Random Forest (RF)~\cite{breiman2001random}, XGB~\cite{chen2016xgboost}, LGBM~\cite{ke2017lightgbm}) and DL architectures, such as LSTM~\cite{hochreiter1997long}, one-dimensional convolutional neural networks (1D-CNN), and a convolutional neural network (CNN)-LSTM hybrid. Finally, it also implements an event-driven, multi-node simulation for scheduling.

We want to measure the impact of using ML-based runtime estimation as part of the scheduling algorithm rather than just reporting the accuracy of such predictions. By simulating controlled scenarios, we can show how much time will be reduced from job wait times by using ML-based runtime estimation compared to other baseline models that do not consider job duration at all (FIFO), or use only resource size, such as Smallest Resource First (SRF).

The majority of the effort that goes into this work will be focused on offline prediction. Models can be created using historical job data and then used to predict runtimes for submitted jobs. There will still be room for online retraining and preemptive scheduling in the future. The experimental analysis has been done solely with the public Alibaba PAI trace~\cite{weng2022mlaas}. It represents an ideal real-world workload size and variety for training, validating, and simulating all aspects of the workflow pipeline under identical conditions.

\needspace{5\baselineskip}
\section{Research Questions}

This thesis addresses the following research questions:
\begin{enumerate}
    \item \textit{Can ML and DL be successfully utilized to solve the job scheduling problem for cloud computing GPU clusters?} Traditional job schedulers use statically defined heuristics that lack an awareness of workload features. This question posits whether data-driven runtime estimates based upon prior knowledge about the characteristics of each job can produce better scheduling decisions with reduced waiting time within a reasonable simulated environment.

    \item \textit{Which predictive modeling paradigm provides the best approach for forecasting the extremely heterogeneous, heavy-tailed nature of runtime in the Alibaba GPU cluster trace?} GPU workloads show extreme heterogeneity in runtime: most jobs complete in less than one minute, whereas many others take hours, days, or weeks. This question compares tree-based ensembles to sequential DL architectures on these extreme distributions to determine which paradigm represents them most accurately.

    \item \textit{How much worse off are resource-aware heuristics that ignore job duration when compared to ML-aided scheduling?} Heuristics like SRF consider the immediate resource demands of a job at the moment the scheduler makes its decision. However, they completely disregard the amount of time a particular job will actually take to execute. This question calculates the performance difference between heuristics and prediction-based scheduling, providing a quantitative measurement of the added benefit provided by predicting execution time.

    \item \textit{What is the empirical performance difference between tree-based ensemble models (XGB, LGBM) and sequential DL architectures (LSTM, 1D-CNN) on tabular data describing GPU cluster workloads?} Recent studies have shown that gradient boosting models typically either match or exceed the performance of DL networks on tabular data. This question examines this hypothesis using large-scale data collected from a real-world production GPU cluster trace, and offers suggestions for model selection to other practitioners working on similar problems.
\end{enumerate}

\needspace{5\baselineskip}
\section{Aims and Objectives}

The specific aims of this research are:
\begin{enumerate}
    \item To process the Alibaba PAI GPU cluster trace and create a feature set that will measure dynamic workload behavior, time submissions for jobs, and current utilization of the cluster as it evolves over time. This utilization is calculated using a sweep-line algorithm to track all jobs running concurrently at every point in time when they were submitted.
    \item To develop and compare the results from training and testing six different ML models under two separate paradigms: tree-based ensemble models (RF, XGB, LGBM) and sequential DL architectures (LSTM, 1D-CNN, CNN-LSTM hybrid). All models used the exact same chronologically ordered data split and were tested/evaluated under exactly the same experimental conditions.
    \item To measure the performance of the ML models based on common regression metrics (mean absolute error (MAE), root mean squared error (RMSE), coefficient of determination (R\textsuperscript{2}), and median absolute error (MedAE)) and to analyze how the models work by evaluating their feature importance (mean decrease in impurity (MDI)) measures.
    \item To build and run a multi-node, event-driven scheduler simulation which will enforce multiple simultaneous CPU, GPU, and memory resources across the nodes within a heterogeneous cluster. The simulation will support four different scheduling algorithms: FIFO, Smallest Resource First (SRF), ML-assisted SJF (SJF-Pred), and an SJF-Oracle policy.
    \item To determine the extent to which ML-assisted SJF reduces average job wait time compared to both FIFO and SRF algorithms and to provide an upper bound on possible improvement through the use of an SJF-Oracle algorithm.
\end{enumerate}

\needspace{5\baselineskip}
\section{Contributions}

This thesis makes the following contributions to the field of cluster resource management:
\begin{enumerate}
    \item \textbf{End-to-end prediction-to-scheduling pipeline.} Raw Alibaba PAI trace ingestion, feature engineering, model training, and event-driven scheduling simulation have been combined into one reproducible pipeline. The pipeline will link prediction accuracy to scheduling results.

    \item \textbf{$O(N \log N)$ sweep-line cluster utilization feature.} Using the sweep-line method, it calculates the number of jobs that were executing when a new job was submitted. It has proven to be one of the top predictive metrics used by many different models using tree-based methods.

    \item \textbf{Decision-Focused Learning framing for cluster scheduling.} This thesis uses a DFL approach rather than a traditional PtO framework. It provides evidence that the primary value of a prediction model is not to minimize statistical error (e.g., MAE); rather, the model becomes truly useful when its predictions enable a scheduler to make better ordinal ranking decisions than it would otherwise make.

    \item \textbf{Six-model cross-paradigm comparison.} Three tree-based ensemble algorithms (RF, XGB, LGBM) and three DL architectures (LSTM, 1D-CNN, CNN-LSTM hybrid) are compared and tested using the same input dataset. Each DL architecture is also trained and evaluated using two separate approaches: categorical and real-valued.

    \item \textbf{Heterogeneous event-driven scheduling simulator.} A discrete-event simulator (DES) is evaluated against several different policies, including FIFO, SRF, ML-assisted SJF (SJF-Pred), and an SJF-Oracle upper-bound policy. The simulator processed over 16,437 test jobs in under 17 seconds per policy.

    \clearpage
    \item \textbf{Quantified scheduling gains across two cluster scales.} By using our ML algorithms, the time required to complete a job was decreased by approximately \textbf{57\%} across both cluster environments from what would have been possible with the baseline FIFO scheduling algorithm. We tested how well each model could approximate the Oracle, which is essentially a hypothetical system that has access to every job's runtime in advance. The results are as follows: the traditional tree-based models (XGB) performed very effectively in the smaller cluster environment; in fact, they approximated \textbf{68.9\%} of the maximum theoretical performance achieved by the Oracle. On the other hand, the DL models (LSTM) performed significantly better on the larger-scale cluster, approximating \textbf{80.4\%} of the maximum theoretical performance achievable by the Oracle.

    \item \textbf{Industrial applicability.} The ML-SJF system was built as a drop-in runtime estimation module that could be used with any SJF-capable scheduler. Since the ML-SJF module replaces inaccurate wall-time provided by users with accurate data-driven ML predictions, this solution works well in combination with other GPU cluster managers (e.g., Kubernetes with GPU plugin), HPC schedulers (e.g., SLURM backfill algorithm), and cloud batch queuing systems (e.g., Yet Another Resource Negotiator (YARN) Capacity Scheduler) so long as there is a trained regression model available along with access to submission-time job metadata.
\end{enumerate}

\needspace{5\baselineskip}
\section{Organization}

This section outlines how the remainder of this thesis is structured. In Chapter~\ref{chp:relatedwork}, we provide the necessary background on GPU cluster management, job scheduling, and the ML and DL methods used in this work, alongside a survey of the related literature. Then, in Chapter~\ref{chp:dataset}, we introduce the Alibaba GPU trace, detail the data preprocessing steps, and present the workload characterization. Next, we present the runtime prediction models in Chapter~\ref{chp:prediction}. Following that, we describe the scheduling simulator framework in Chapter~\ref{chp:simulation}. Finally, we report the experimental results in Chapter~\ref{chp:results}, and conclude by summarizing the findings and proposing future directions for research in Chapter~\ref{chp:conclusion}.